% 本文件由 LaTeX 排版 Agent 根据有效 translation.json 自动生成。
% author=Miscellaneous; work_id=0852; title=关于阿布加王与阿戴宗徒之各书摘录

\chapter{关于\Person{阿布加}王与\Person{阿戴}宗徒之各书摘录}

% structure: p0001 source=h1 target=omitted reason=page-title-already-rendered-by-parent

% structure: p0002 source=h2 target=section reason=relative-heading-rank; base=h2

\section{1}

\Emph{关于真福宗徒\Person{阿戴}——取自他在\Place{厄德撒}城于\Person{阿布加}王及全城集会前所作之教导。}\par

及至祂进入坟墓，便复活了，与许多人一同从坟墓出来。那些看守坟墓的人并未看见祂如何从坟墓出来；但天上的守望者——他们是祂复活的宣告者和报信者。因为祂若不愿意，就不会死去，因祂是死亡的主，即离开\Emph{此生}的门；若祂不喜欢，也不会穿上身体，因祂本身就是身体的创造者。因为那引导祂屈尊由童贞女降生的旨意，也同样使祂进一步下降至受死亡的痛苦。——\Emph{稍后}（\Emph{我们读到}）：因为，尽管祂的外表与人相似，但祂的能力、知识和权柄却是天主的。\par

% structure: p0005 source=h2 target=section reason=relative-heading-rank; base=h2

\section{2}

\Emph{取自宗徒\Person{阿戴}在\Place{厄德撒}城所述之教导。}\par

你们知道我曾对你们说过：从人身体出来的灵魂，没有一个处于\Emph{死亡的控制}之下，而都是活着并继续存在，并且为他们预备了居所和安息之所。因为灵魂的理性能力不会停止，知识也不会消失，因为它是永生天主的肖像。因为它不像肉体那样没有知觉——肉体对于已经掌控它的腐败没有知觉。然而，报应和赏报必须等到它结合其形体才能获得，因为所经历的不只属于灵魂本身，也属于它曾暂时居住的形体。但那些悖逆者，即不认识天主的人，那时将徒然后悔。\par

% structure: p0008 source=h2 target=section reason=relative-heading-rank; base=h2

\section{3}

\Emph{取自宗徒\Person{阿戴}在\Place{厄德撒}城所述之书信。}\par

留意你们所领受的这项职务，并怀着恐惧和战栗坚守其中，每日尽忠职守。不可因疏忽的习惯而事奉，而要以信德的谨慎行事。愿基督的赞美不绝于你们口中，愿在祈祷之时没有任何倦怠之感。留意你们所持守的真理，以及你们所接受的真理教导，还有我交付给你们的救恩教导。因为在基督的审判台前，这些将要向你们追讨，当祂与牧者及监督算账时，当祂从商人那里收回祂的银子，连同他们所教导的利息时。因为祂是王的儿子，要去领受王国，并且祂要回来，使众人复活得生命。\par

% structure: p0011 source=h2 target=section reason=relative-heading-rank; base=h2

\section{4}

\Person{阿戴}在\Place{厄德撒}和\Place{美索不达米亚}宣讲（他来自\Place{帕内乌斯}），当时正值\Person{阿布加}王在位。当他身处\Place{佐费尼人}之中时，\Person{阿布加}的儿子\Person{塞韦鲁斯}派人将他杀害于\Place{阿格勒·哈斯纳}，也杀了一位与他同行的年轻门徒。\par

% structure: p0013 source=h2 target=section reason=relative-heading-rank; base=h2

\section{5}

71. 与\Person{纳齐盎}。因为他们不容许那七十二人的选召有所短缺，同样也不容许那十二人的选召短缺。这\Emph{人}属于那七十二人：或许他是宗徒\Person{阿戴}的门徒。\par

% structure: p0015 source=h2 target=section reason=relative-heading-rank; base=h2

\section{6}

\Emph{取自\Person{圣母玛利亚}离世，以及我主\Person{耶稣基督}的诞生与童年。第二卷。}\par

在第三百四十五年，后\Place{提史林}月，\Person{圣母玛利亚}从家中出来，前往\Person{基督}的坟墓；因为她每日都去那里哀哭。但犹太人，在\Person{基督}死后立即夺取了坟墓，并在墓门口堆上大石头。他们在坟墓和\Place{哥耳哥达}派了守卫，命令他们：若有人去坟墓或\Place{哥耳哥达}祈祷，应立即处死。犹太人拿走了我们主的十字架，以及另外两个十字架，还有刺伤我们救主的长枪，钉入祂手足的钉子，以及祂所穿的戏弄袍子；并将它们藏起来，免得——如他们所说——有君王或首领前来查问\Person{基督}被处死之事。\par

守卫进去对司铎说：\Person{玛利亚}傍晚和早晨来此祈祷。\Place{耶路撒冷}因\Person{圣母玛利亚}而骚动。司铎去见审判官，对他说：我主，请下令禁止\Person{玛利亚}到坟墓和\Place{哥耳哥达}祈祷。正当他们商议时，看，信函从\Place{厄德撒}城的\Person{阿布加}王送到了由\Person{提庇留}皇帝任命的检察官\Person{萨比纳}那里，检察官\Person{萨比纳}的权柄直达\Place{幼发拉底河}。因为宗徒\Person{阿戴}，七十二宗徒之一，曾前往\Place{厄德撒}建造了一所教会，并治愈了\Person{阿布加}王所患的疾病——因为\Person{阿布加}王爱\Person{耶稣基督}，并时常打听祂的事；当\Person{基督}被处死，\Person{阿布加}王听说犹太人将祂钉死在十字架上，便大为不悦；于是\Person{阿布加}起身骑马，直抵\Place{幼发拉底河}，因为他想上\Place{耶路撒冷}去将其摧毁；当\Person{阿布加}来到\Place{幼发拉底河}时，他心中思忖：我若过河，便会与\Person{提庇留}皇帝结下仇怨。于是\Person{阿布加}写了几封信，送给检察官\Person{萨比纳}，\Person{萨比纳}将它们呈给\Person{提庇留}皇帝。\Person{阿布加}给\Person{提庇留}皇帝写信如下：——\par

\Quote{从\Place{厄德撒}城王\Person{阿布加}，致陛下我主\Person{提庇留}，大安！为使陛下不致因我而不悦，我未渡过\Place{幼发拉底河}；因我一直想上\Place{耶路撒冷}将其摧毁，因她杀害了\Person{基督}，一位精妙的医治者。但陛下您作为一位统辖全地和我们的大君主，请派兵为我向\Place{耶路撒冷}居民施行审判。因为应让陛下知道，我盼望您能为我对那些钉死祂的人施行审判。}\par

撒彼纳收信后，将其呈给皇帝提庇留。提庇留皇帝读后大怒，欲毁灭并杀尽所有犹太人。耶路撒冷的居民听闻此事，惊慌失措。司祭们去见总督，对他说：\Quote{我主，请派人吩咐玛利亚，不可到圣墓和哥耳哥达祈祷。}判官对司祭说：\Quote{你们自己去，随意吩咐和警告她吧。}\par

% structure: p0021 source=h2 target=section reason=relative-heading-rank; base=h2

\section{7}

\Emph{选自圣雅各伯导师论偶像堕落的讲道。}\par

他前往厄德萨，在那里发现了一项伟大的工程：\par

国王成了教会的工人，正在建造它。\par

宗徒阿戴站立其中，如同建筑师，\par

阿布加尔王摘下王冠，与他一同建造。\par

当宗徒与国王同心合意时，\par

有何偶像不在他们面前倾倒？\par

撒旦从门徒面前逃往巴比伦地，\par

而十字架的故事已先他传到加色丁人之地。\par

他说，当他们嘲弄黄道十二宫的征兆时，他什么也不是。\par

% structure: p0032 source=h2 target=section reason=relative-heading-rank; base=h2

\section{8}

\Emph{选自关于安提约基雅城的讲道。}\par

西满分得罗马，若望分得厄弗所，多默分得印度，阿戴分得亚述地区。他们被分别派往各自分得的区域后，专心致力于带领\Emph{若干}地区成为门徒。\par

\SourceRecord{Translated by B.P. Pratten.；Ante-Nicene Fathers；Vol. 8.；Edited by Alexander Roberts, James Donaldson, and A. Cleveland Coxe.；Buffalo, NY: Christian Literature Publishing Co.,；1886.；<http://www.newadvent.org/fathers/0852.htm>.}
